\slajdid{04-s039}
\begin{frame}[label=04-s039]{Prikaz sa pokretnom tačkom}
% element: 04-s039:e01 — Prethodni prikaz sa fiksnom tačkom i floating point
% element: 04-s039:e02 — Dva registra ili dva dela jednog registra: mantisa i eksponent
% element: 04-s039:e03 — Formula V=mantisa×osnova^eksponent
% element: 04-s039:e04 — Binarna mantisa kao znak/apsolutna vrednost i fiksna tačka iza MSB-a
% element: 04-s039:e05 — Eksponent kao označeni ceo broj sa ofsetom; veći opsezi i tačnost
Do sada: \alert{fiksna tačka}. Postoji i prikaz sa \alert{pokretnom tačkom} — \textit{floating point number}.\par
Broj ima dva registra ili jedan registar sa dva dela: mantisa i eksponent.
\[\alert{V=\mathit{mantisa}\times\mathit{osnova}^{\mathit{eksponent}}}\]
\begin{columns}[T,onlytextwidth]
\column{.475\textwidth}\NaslovBloka{Mantisa\\u binarnom sistemu}
Označeni binarni broj: znak i apsolutna vrednost, sa fiksnom tačkom iza MSB-a.
\column{.475\textwidth}\NaslovBloka{Eksponent\\u binarnom sistemu}
Označeni celobrojni binarni broj, sa ofsetom definisanim standardom.
\end{columns}
\par\vspace{4mm}\Rezultat{Dozvoljava veće opsege i tačnost prikazivanja.}
\end{frame}
